\documentclass[11pt]{article}
% The preamble every project's paper shares: % The preamble every project's paper shares: % The preamble every project's paper shares: \input{../../preamble} from
% projects/<project>/paper/. Boards come from chessboard and xskak, plots from
% pgfplots; every number a paper prints comes from its generated.tex.
\usepackage[LSBC4,T1]{fontenc}
\usepackage[utf8]{inputenc}
\usepackage{lmodern}
\usepackage[margin=1in]{geometry}
\usepackage{microtype}
\usepackage{amsmath}
\usepackage{booktabs}
\usepackage{tabularx}
\usepackage{array}
\usepackage{xcolor}
\usepackage[obeyspaces,spaces]{url}
\usepackage{xskak}
\usepackage{pgfplots}
\pgfplotsset{compat=1.18}
\usepackage{caption}
\usepackage{subcaption}
\usepackage{enumitem}
\usepackage[section]{placeins}
\usepackage[hidelinks]{hyperref}

\setchessboard{boardfontencoding=LSBC4, boardfontfamily=skaknew, blackfieldcolor=black!22,
  whitefieldcolor=white, setfontcolors, showmover=false, label=true, labelfontsize=6pt,
  boardfontsize=14pt}

% A move or line in UCI, and a position in FEN, both breakable at a slash.
\newcommand{\uci}[1]{\texttt{#1}}
\newcommand{\fen}[1]{{\footnotesize\path{#1}}}
\newcolumntype{L}{>{\raggedright\arraybackslash}X}

\setlength{\parskip}{0.4em}
\setlength{\parindent}{0pt}
 from
% projects/<project>/paper/. Boards come from chessboard and xskak, plots from
% pgfplots; every number a paper prints comes from its generated.tex.
\usepackage[LSBC4,T1]{fontenc}
\usepackage[utf8]{inputenc}
\usepackage{lmodern}
\usepackage[margin=1in]{geometry}
\usepackage{microtype}
\usepackage{amsmath}
\usepackage{booktabs}
\usepackage{tabularx}
\usepackage{array}
\usepackage{xcolor}
\usepackage[obeyspaces,spaces]{url}
\usepackage{xskak}
\usepackage{pgfplots}
\pgfplotsset{compat=1.18}
\usepackage{caption}
\usepackage{subcaption}
\usepackage{enumitem}
\usepackage[section]{placeins}
\usepackage[hidelinks]{hyperref}

\setchessboard{boardfontencoding=LSBC4, boardfontfamily=skaknew, blackfieldcolor=black!22,
  whitefieldcolor=white, setfontcolors, showmover=false, label=true, labelfontsize=6pt,
  boardfontsize=14pt}

% A move or line in UCI, and a position in FEN, both breakable at a slash.
\newcommand{\uci}[1]{\texttt{#1}}
\newcommand{\fen}[1]{{\footnotesize\path{#1}}}
\newcolumntype{L}{>{\raggedright\arraybackslash}X}

\setlength{\parskip}{0.4em}
\setlength{\parindent}{0pt}
 from
% projects/<project>/paper/. Boards come from chessboard and xskak, plots from
% pgfplots; every number a paper prints comes from its generated.tex.
\usepackage[LSBC4,T1]{fontenc}
\usepackage[utf8]{inputenc}
\usepackage{lmodern}
\usepackage[margin=1in]{geometry}
\usepackage{microtype}
\usepackage{amsmath}
\usepackage{booktabs}
\usepackage{tabularx}
\usepackage{array}
\usepackage{xcolor}
\usepackage[obeyspaces,spaces]{url}
\usepackage{xskak}
\usepackage{pgfplots}
\pgfplotsset{compat=1.18}
\usepackage{caption}
\usepackage{subcaption}
\usepackage{enumitem}
\usepackage[section]{placeins}
\usepackage[hidelinks]{hyperref}

\setchessboard{boardfontencoding=LSBC4, boardfontfamily=skaknew, blackfieldcolor=black!22,
  whitefieldcolor=white, setfontcolors, showmover=false, label=true, labelfontsize=6pt,
  boardfontsize=14pt}

% A move or line in UCI, and a position in FEN, both breakable at a slash.
\newcommand{\uci}[1]{\texttt{#1}}
\newcommand{\fen}[1]{{\footnotesize\path{#1}}}
\newcolumntype{L}{>{\raggedright\arraybackslash}X}

\setlength{\parskip}{0.4em}
\setlength{\parindent}{0pt}

% Written by projects/maia-calibration/paper/generate.mjs; do not edit by hand.

\newcommand{\AppCommit}{\texttt{ad7fa98}}
\newcommand{\Draws}{1{,}000}
\newcommand{\Seed}{7}
\newcommand{\MinFitMoves}{300}
\newcommand{\LichessMoves}{25{,}000}
\newcommand{\LichessGames}{13{,}785}
\newcommand{\FitMoves}{12{,}474}
\newcommand{\HeldMoves}{12{,}526}
\newcommand{\HisMoves}{33{,}522}
\newcommand{\HisGames}{640}
\newcommand{\KeptOffset}{250}
\newcommand{\PunishExact}{521}
\newcommand{\PunishAll}{684}
\newcommand{\PunishGames}{60}
\newcommand{\BenchMonth}{June 2025}
\newcommand{\BenchPerBand}{5{,}000}
\newcommand{\BenchPerGame}{8}
\newcommand{\BenchSeen}{150{,}181}
\newcommand{\BenchMinPly}{11}
\newcommand{\BenchMinClock}{30}
\newcommand{\BenchGapMin}{3.0}
\newcommand{\BenchGapMax}{3.9}
\newcommand{\BenchTop}{51.5}
\newcommand{\BenchConf}{54.9}
\newcommand{\FirstChoice}{12{,}875}
\newcommand{\TestTable}{%
\begin{tabular}{lrrrl}
\toprule
test & raw & tempered & tempered $-$ raw & \\
\midrule
1. Lichess, held-out half (log-likelihood per move) & $-$1.4944 & $-$1.4915 & 0.0029 [0.0014, 0.0044] & passes \\
2. the player's games (log-likelihood per move) & $-$1.3959 & $-$1.4026 & $-$0.0068 [$-$0.0079, $-$0.0057] & fails \\
3. punish events, verified (Brier score) & 0.1527 & 0.1542 & 0.0015 [$-$0.0005, 0.0037] & passes \\
\phantom{3.} punish events, all (Brier score) & 0.1550 & 0.1563 & 0.0013 [$-$0.0005, 0.0032] & \\
\bottomrule
\end{tabular}
}
\newcommand{\TestOneDiff}{0.0029}
\newcommand{\TestOneCI}{[0.0014, 0.0044]}
\newcommand{\TestTwoLoss}{0.0068}
\newcommand{\TestTwoLossCI}{[0.0057, 0.0079]}
\newcommand{\TestThreeDiff}{0.0015}
\newcommand{\TestThreeCI}{[$-$0.0005, 0.0037]}
\newcommand{\TestThreeAllDiff}{0.0013}
\newcommand{\TestThreeAllCI}{[$-$0.0005, 0.0032]}
\newcommand{\PunishRaw}{0.153}
\newcommand{\PunishTempered}{0.154}
\newcommand{\TempMin}{1.02}
\newcommand{\TempMax}{1.13}
\newcommand{\FitMovesMin}{620}
\newcommand{\TempFigure}{%
\begin{tikzpicture}
\begin{axis}[ybar, bar width=9pt, width=0.85\textwidth, height=0.36\textwidth, ymin=0.95, ymax=1.16,
  xmin=-0.6, xmax=10.6, xtick={0,1,2,3,4,5,6,7,8,9,10}, xticklabels={{<1100},{1100},{1200},{1300},{1400},{1500},{1600},{1700},{1800},{1900},{2000+}}, xticklabel style={font=\scriptsize},
  ylabel={temperature $T$}, xlabel={Maia-2 band (Lichess rating)}, ymajorgrids, enlarge x limits=false,
  y tick label style={/pgf/number format/fixed, /pgf/number format/precision=2}]
\addplot[fill=black!35, draw=black] coordinates {(0,1.0822) (1,1.0639) (2,1.0847) (3,1.0878) (4,1.1254) (5,1.0583) (6,1.0682) (7,1.0179) (8,1.0339) (9,1.0709) (10,1.0793)};
\draw[dashed, red!70!black] (axis cs:-0.6,1) -- (axis cs:10.6,1);
\end{axis}
\end{tikzpicture}
}
\newcommand{\LichGapMin}{2.0}
\newcommand{\LichGapMax}{5.2}
\newcommand{\LichTGapMin}{$-$1.1}
\newcommand{\LichTGapMax}{3.4}
\newcommand{\HisGapMin}{$-$1.3}
\newcommand{\HisGapMax}{1.1}
\newcommand{\HisTGapMin}{1.8}
\newcommand{\HisTGapMax}{4.6}
\newcommand{\HisBandLow}{<1100}
\newcommand{\HisBandHigh}{1400}
\newcommand{\GapFigure}{%
\begin{tikzpicture}
\begin{axis}[name=lichess, ylabel={confidence $-$ played, points}, width=0.48\textwidth, height=0.4\textwidth, xmin=-0.6, xmax=10.6, ymin=-6, ymax=6,
  xtick={0,1,2,3,4,5,6,7,8,9,10}, xticklabels={{<1100},{1100},{1200},{1300},{1400},{1500},{1600},{1700},{1800},{1900},{2000+}}, xticklabel style={font=\tiny, rotate=60, anchor=east}, grid=major,
  title={Lichess, held-out half}, legend pos=south west, legend cell align=left, legend style={font=\scriptsize}]
\addplot[black, mark=*, mark size=1.2pt] coordinates {(0,3.68) (1,5.00) (2,2.85) (3,2.70) (4,2.90) (5,5.22) (6,1.95) (7,2.41) (8,3.56) (9,4.09) (10,3.51)};
\addlegendentry{raw}
\addplot[red!70!black, dashed, mark=square*, mark size=1.2pt] coordinates {(0,0.85) (1,2.77) (2,0.13) (3,-0.18) (4,-1.12) (5,3.35) (6,-0.28) (7,1.81) (8,2.43) (9,1.77) (10,0.98)};
\addlegendentry{tempered}
\draw[gray] (axis cs:-0.6,0) -- (axis cs:10.6,0);
\end{axis}
\begin{axis}[at={(lichess.east)}, anchor=west, xshift=0.8cm, width=0.48\textwidth, height=0.4\textwidth, xmin=-0.6, xmax=10.6, ymin=-6, ymax=6,
  xtick={0,1,2,3,4,5,6,7,8,9,10}, xticklabels={{<1100},{1100},{1200},{1300},{1400},{1500},{1600},{1700},{1800},{1900},{2000+}}, xticklabel style={font=\tiny, rotate=60, anchor=east}, grid=major,
  title={the player's games}, legend pos=south west, legend cell align=left, legend style={font=\scriptsize}]
\addplot[black, mark=*, mark size=1.2pt] coordinates {(0,-1.30) (1,0.23) (2,1.08) (3,0.55) (4,-0.64)};
\addlegendentry{raw}
\addplot[red!70!black, dashed, mark=square*, mark size=1.2pt] coordinates {(0,-4.04) (1,-1.89) (2,-1.75) (3,-2.30) (4,-4.62)};
\addlegendentry{tempered}
\draw[gray] (axis cs:-0.6,0) -- (axis cs:10.6,0);
\end{axis}
\end{tikzpicture}
}
\newcommand{\BandTable}{%
\begin{tabular}{lrrrrrrrrr}
\toprule
 & & & \multicolumn{3}{c}{Lichess, held-out half} & \multicolumn{4}{c}{the player's games} \\
\cmidrule(lr){4-6}\cmidrule(lr){7-10}
band & fitting moves & $T$ & played & raw & tempered & moves & played & raw & tempered \\
\midrule
<1100 & 2{,}537 & 1.082 & 46.2 & 49.8 & 47.0 & 6{,}618 & 53.0 & 51.7 & 48.9 \\
1100 & 729 & 1.064 & 49.2 & 54.2 & 52.0 & 5{,}906 & 54.1 & 54.3 & 52.2 \\
1200 & 833 & 1.085 & 51.3 & 54.2 & 51.5 & 6{,}893 & 53.1 & 54.1 & 51.3 \\
1300 & 946 & 1.088 & 52.2 & 54.9 & 52.0 & 13{,}222 & 54.6 & 55.2 & 52.3 \\
1400 & 814 & 1.125 & 53.1 & 56.0 & 52.0 & 883 & 56.5 & 55.9 & 51.9 \\
1500 & 823 & 1.058 & 49.2 & 54.4 & 52.5 &  & & & \\
1600 & 841 & 1.068 & 52.8 & 54.7 & 52.5 &  & & & \\
1700 & 1{,}001 & 1.018 & 55.2 & 57.6 & 57.0 &  & & & \\
1800 & 833 & 1.034 & 53.4 & 56.9 & 55.8 &  & & & \\
1900 & 620 & 1.071 & 53.3 & 57.4 & 55.1 &  & & & \\
2000+ & 2{,}497 & 1.079 & 54.2 & 57.7 & 55.2 &  & & & \\
\bottomrule
\end{tabular}
}


\title{Is Maia-2 too sure of itself?\\[0.3em]
\large A temperature for each rating band, fitted on Lichess and tested on the player's games}
\author{Chess Coach research lab}
\date{}

\begin{document}
\maketitle

\begin{abstract}
The coaching app reads Maia-2's probabilities as chances: the odds that a mistake is punished,
how findable a move is, the moves its Maia bot samples. The lab's benchmark found Maia-2 more
confident than right on Lichess in every band. We fitted a temperature for each band on
\FitMoves{} Lichess moves and tested it by a rule fixed before any number was computed. It
improves the held-out Lichess moves and does not hurt the punish odds, but it makes the player's
own games less likely, so it is not adopted. On his games, at the rating offset the app uses,
Maia-2 is already about as sure as it is right, and softening it makes it less sure than right.
The app keeps the raw policy. Scored by the app's own Maia-2, the Lichess sample also reproduces
the benchmark's first-choice count exactly.
\end{abstract}

\section{The question}

The coaching app reads Maia-2's probabilities as chances. A review says how often players at the
opponent's level find the reply that punishes a mistake; a drill says how findable a move is at
the player's level; the Maia bot samples its moves from them. The lab's benchmark of Maia-2, on
\LichessMoves{} moves from rated rapid games on Lichess in \BenchMonth, after Maia-2's training
window, found it more confident than right in every band: its first choice is the move played
\BenchTop\% of the time, while the mean probability it gives that choice is \BenchConf\%, and
the gap runs from \BenchGapMin{} to \BenchGapMax{} points across the benchmark's bands.

A common remedy is a temperature: each probability $p$ becomes $p^{1/T}$, divided by the same
sum over the position's legal moves. A temperature above one softens the policy and below one
sharpens it; either way the order of the moves stays, so the top move and every ranking the app
reads stay as they are. Does a temperature for each band, fitted on Lichess, make the
probabilities truer on the player's own games?

\section{The data}

\paragraph{Fitting.} The benchmark's own sample, drawn again by its sampler, unchanged:
\BenchPerBand{} moves per band of the benchmark from \BenchSeen{} streamed games, at most
\BenchPerGame{} per player and game, from ply \BenchMinPly{} on with both clocks at
\BenchMinClock{} seconds or more. Each move is scored by the Maia-2 the app ships, through its engine service in Node, with
each side at its own Lichess rating's band. The \LichessGames{} games, ordered by their ids,
alternate between a fitting half (\FitMoves{} moves) and a held-out half (\HeldMoves{} moves).

\paragraph{Testing.} The player's games, never used to fit: the \HisMoves{} moves of the
rating-offset project in \HisGames{} games, with both players' chess.com ratings converted at the
offset that project keeps (\KeptOffset); and the punish events of the punish-odds project, each
priced as the app prices it, at the replier's band facing an equal opponent. The punish events'
primary set holds the \PunishExact{} replies whose evidence the review's deeper pass verified;
all \PunishAll{} non-book replies, in \PunishGames{} games, are reported beside it.

\paragraph{Parity.} Scored by the app's Maia-2, the move played is its first choice for
\FirstChoice{} of the \LichessMoves{} moves, exactly the count the benchmark's Maia-2, the
reference implementation, gives on the same moves.

\section{The candidate and the rule}

Each band's temperature maximises the likelihood of the moves played by movers in that band in
the fitting half; a band with fewer than \MinFitMoves{} fitting moves would keep $T = 1$. The rule
was committed before any number was computed: the temperatures are adopted only if all three
tests hold, each judged by the 95\% interval of the difference over \Draws{} bootstrap draws of
games.

\begin{enumerate}[nosep]
\item On the Lichess held-out half, the mean log-likelihood per move of the moves played rises.
\item On the player's games, the same.
\item On the punish events, the Brier score of the punish odds, against whether the punishing
reply was played, does not rise.
\end{enumerate}

\section{Results}

Not adopted: the temperatures pass the first and third tests and fail the second
(Table~\ref{tab:tests}). Fitted on Lichess, every band's temperature is above one, from
\TempMin{} to \TempMax{} (Figure~\ref{fig:temps}), so every band is softened; each had at least
\FitMovesMin{} fitting moves. On the Lichess held-out half, the softer policy raises the mean
log-likelihood of the moves played by \TestOneDiff{} per move \TestOneCI. On the player's games
it lowers it by \TestTwoLoss{} per move \TestTwoLossCI, and the rule stops there. On the punish
events the Brier score moves from \PunishRaw{} to \PunishTempered, a rise of \TestThreeDiff{}
\TestThreeCI{} whose interval holds zero, so the third test passes.

\begin{table}[!htbp]
\centering\scriptsize
\setlength{\tabcolsep}{4pt}
\TestTable
\caption{The three tests, raw against tempered, with the 95\% interval of the difference over
\Draws{} bootstrap draws of games. A higher log-likelihood and a lower Brier score are better.
All non-book punish events are reported beside the verified ones and decide nothing.}
\label{tab:tests}
\end{table}

\begin{figure}[!htbp]
\centering
\TempFigure
\caption{Each band's temperature, fitted on the Lichess fitting half. Above the dashed line a
band's policy is softened.}
\label{fig:temps}
\end{figure}

\paragraph{Why it does not transfer.} Figure~\ref{fig:gaps} and Table~\ref{tab:bands} compare
Maia-2's mean probability of its first choice with how often that choice is the move played. On
the Lichess held-out half Maia-2 is more confident than right in every band, by \LichGapMin{} to
\LichGapMax{} points, and the temperature brings that to between \LichTGapMin{} and
\LichTGapMax. On the player's games there is little to correct: in the bands he and his opponents
occupy (\HisBandLow{} to \HisBandHigh), the gap runs from \HisGapMin{} to \HisGapMax{} points,
smaller in every band than the smallest on Lichess, and softening turns it into
underconfidence of \HisTGapMin{} to \HisTGapMax{} points. One reading, not tested here: the
rating offset was chosen on these games for the same likelihood, and a higher band is a sharper
policy, so the offset already sets how sure Maia-2 is about them.

\begin{figure}[!htbp]
\centering
\GapFigure
\caption{Maia-2's mean probability of its first choice less the share of moves where that choice
was played, in points, by band: above zero it is more confident than right.}
\label{fig:gaps}
\end{figure}

\begin{table}[!htbp]
\centering\scriptsize
\setlength{\tabcolsep}{4pt}
\BandTable
\caption{By band: the fitting moves and the fitted temperature, then, on the Lichess held-out
half and on the player's games, the share of moves where Maia-2's first choice was played and its
mean probability, raw and tempered (percent).}
\label{tab:bands}
\end{table}

\section{What it means for the app}

Nothing to port. On this player's games Maia-2, at the bands the app's offset gives, is about
as sure as it is right, so the punish odds, findability and the Maia bot keep the raw policy.
The overconfidence the benchmark measured is real on Lichess; it is not what this player's games
show.

\section{Limitations}

\begin{itemize}[nosep]
\item One temperature per band corrects confidence on average. It cannot fix a policy that
ranks the moves wrongly, and it says nothing about positions where Maia-2 is sure for a reason.
\item The fitting sample is one month of Lichess rapid. The player's games are chess.com rapid,
through a rating conversion that is itself fitted.
\end{itemize}

\end{document}
